\documentclass[12pt, a4paper]{article}
\usepackage[a4paper, top=1.2cm, bottom=1.5cm, left=1.5cm, right=1.5cm]{geometry}
\usepackage{amsmath, amssymb}
\usepackage{fontspec}
\usepackage[bidi=basic]{babel}
\babelprovide[main, import, onchar=ids fonts]{arabic}
\babelprovide[import, onchar=ids fonts]{english}
\babelfont{rm}{Noto Serif}
\babelfont[arabic]{rm}{Amiri} % الخط الأميري
\usepackage{enumitem}
\usepackage{tabularx}
\usepackage{array}
\usepackage{graphicx}
\usepackage{tikz}
\usetikzlibrary{calc, angles, quotes}
\usepackage{eso-pic}
\usepackage{tcolorbox}

% ------------------------------------------------------------------
% مهم: الصفحة باتجاه عربي (من اليمين لليسار)، وهذا يعكس أماكن الرموز
% والزوايا داخل رسومات TikZ. لذلك نضع كل رسمة داخل صندوق إنجليزي (LTR)
% ------------------------------------------------------------------
\newcommand{\LTRfig}[1]{\mbox{\foreignlanguage{english}{#1}}}

% قطاع دائري بنفس شكل الكتاب: الرأس على اليسار والقوس على اليمين
% #1 = قياس الزاوية ، #2 = نصف القطر في الرسم (cm) ، #3 = طول نصف القطر المكتوب
\newcommand{\sectorfig}[3]{%
\LTRfig{%
\begin{tikzpicture}[thick]
  \pgfmathsetmacro{\s}{-12}          % ميل نصف القطر السفلي
  \pgfmathsetmacro{\e}{\s + #1}      % اتجاه نصف القطر العلوي
  \pgfmathsetmacro{\m}{\s + #1/2}    % منتصف الزاوية
  \draw[violet, fill=violet!4] (0,0) -- (\s:#2)
        arc[start angle=\s, end angle=\e, radius=#2] -- cycle;
  \draw[red] (\s:0.25*#2) arc[start angle=\s, end angle=\e, radius=0.25*#2];
  \node at (\m:0.45*#2) {$#1^\circ$};
  \node[below=3pt] at (\s:0.5*#2) {$#3$};
\end{tikzpicture}}}

\begin{document}

% إطار الصفحات
\AddToShipoutPictureBG{%
  \LTRfig{%
  \begin{tikzpicture}[remember picture, overlay]
    \draw[line width=1.2pt] ([shift={(0.8cm,-0.8cm)}]current page.north west) rectangle ([shift={(-0.8cm,0.8cm)}]current page.south east);
  \end{tikzpicture}}%
}

% الترويسة
\noindent
\begin{minipage}[t]{0.38\textwidth}
\vspace{0pt}
\textbf{المملكة الأردنية الهاشمية} \\
\textbf{وزارة التربية والتعليم} \\
\textbf{مدرسة عنبة الثانوية الشاملة للبنات} \\
\textbf{المبحث: الرياضيات}
\end{minipage}%
\begin{minipage}[t]{0.24\textwidth}
\vspace{0pt}
\centering
\includegraphics[width=2.3cm]{logo.png}
\end{minipage}%
\begin{minipage}[t]{0.38\textwidth}
\vspace{0pt}
\textbf{الصف: العاشر} \\
\textbf{الشعبة: ( \hspace{1.5cm} )} \\
\textbf{اليوم والتاريخ: \hspace{0.5cm} / \hspace{0.5cm} / 202 م} \\
\textbf{الزمن: ساعة واحدة}
\end{minipage}

\vspace{0.3cm}
\hrule height 1.2pt
\vspace{0.08cm}
\hrule height 0.4pt
\vspace{0.5cm}

\begin{center}
\textbf{\Large اختبار الرياضيات}
\end{center}

\vspace{0.2cm}
\noindent \textbf{اسم الطالبة:} \rule{10cm}{0.5pt}

\vspace{0.5cm}
\noindent\textbf{\large السؤال الأول: اختاري رمز الإجابة الصحيحة في كل مما يأتي (18 علامة):}
\vspace{0.2cm}

\begin{enumerate}[label=\textbf{\arabic*-}, itemsep=0.35cm, leftmargin=*]

\item يمكن حل المعادلة الآتية $x^3 + 4x^2 = 5x$ عن طريق:
\vspace{0.2cm} \\
\noindent\begin{tabularx}{\linewidth}{XXXX}
(A) إخراج عامل مشترك & (B) بالتجميع & (C) مجموع مكعبين & (D) الفرق بين مكعبين
\end{tabularx}

\item جذور المعادلة $x^3 + 2x^2 - x - 2 = 0$ هي:
\vspace{0.2cm} \\
\noindent\begin{tabularx}{\linewidth}{XXXX}
(A) جذراً واحداً & (B) $-1, 1, 2$ & (C) $-3, -2, 2$ & (D) $-1, 1, -2$
\end{tabularx}

\item إذا كانت قيمة المميز سالبة فإن عدد حلول النظام هو:
\vspace{0.2cm} \\
\noindent\begin{tabularx}{\linewidth}{XXXX}
(A) حل واحد فقط & (B) لا يوجد حل & (C) حلين للنظام & (D) ثلاث حلول
\end{tabularx}

\item أي الأزواج المرتبة الآتية يمثل حلاً لنظام المعادلات:
$\begin{cases} y = x^2 - 5x + 6 \\ y = -x^2 + 2x + 3 \end{cases}$
\vspace{0.2cm} \\
\noindent\begin{tabularx}{\linewidth}{XXXX}
(A) $(3,0)$ & (B) $(0,3)$ & (C) $(0,2)$ & (D) $(2,0)$
\end{tabularx}

\item القطعة المستقيمة الواصلة بين نقطتين على الدائرة وتمر بالمركز تسمى:
\vspace{0.2cm} \\
\noindent\begin{tabularx}{\linewidth}{XXXX}
(A) القاطع & (B) المماس & (C) القطر & (D) نصف القطر
\end{tabularx}

\item الوتران اللذان يبعدان المسافة نفسها عن المركز هما وتران:
\vspace{0.2cm} \\
\noindent\begin{tabularx}{\linewidth}{XXXX}
(A) متطابقان & (B) متوازيان & (C) متعامدان & (D) $(B+C)$
\end{tabularx}

\item قياس الزاوية المركزية يساوي:
\vspace{0.2cm} \\
\noindent\begin{tabularx}{\linewidth}{XX}
(A) نصف قياس الزاوية المحيطية & (B) ربع قياس الزاوية المحيطية \\[0.2cm]
(C) مثلي قياس الزاوية المحيطية & (D) ثلث قياس الزاوية المحيطية
\end{tabularx}

\item طول القوس بدلالة ($\pi$) في الشكل الآتي: \hfill \raisebox{-0.9cm}{\sectorfig{25}{4}{8}} \hspace{1.5cm}
\vspace{0.2cm} \\
\noindent\begin{tabularx}{\linewidth}{XXXX}
(A) $\dfrac{40\pi}{9}$ & (B) $\dfrac{10\pi}{9}$ & (C) $3.489$ & (D) $13.956$
\end{tabularx}

\item العددان اللذان مجموعهما 12 والفرق بين مربعيهما 24 هما:
\vspace{0.2cm} \\
\noindent\begin{tabularx}{\linewidth}{XXXX}
(A) $(-7,-5)$ & (B) $(7,5)$ & (C) $(8,4)$ & (D) $(-8,-4)$
\end{tabularx}

\end{enumerate}

\newpage
% بداية الصفحة الثانية
\noindent\textbf{\large السؤال الثاني: أحلّ نظام المعادلات الآتي، ثم أتحقق من صحة الحل (10 علامات):}
\vspace{0.3cm}
\begin{align*}
y &= x^2 + 6x - 3 \\
y &= 2x - 3
\end{align*}
\vspace{6cm} % مساحة للحل

\hrule height 0.4pt
\vspace{0.5cm}

\noindent\textbf{\large السؤال الثالث: أجد طول القوس ومساحة القطاع ومحيط القطاع في الشكل المجاور: (10 علامات)}
% (ملاحظة: نفس شكل سؤال 6 صفحة 44)
\begin{center}
\sectorfig{60}{4}{12\text{ cm}}
\end{center}
\vspace{2.5cm} % مساحة للحل

\hrule height 0.4pt
\vspace{0.5cm}

\noindent\textbf{\large السؤال الرابع: مزرعة مستطيلة الشكل محيطها $16\text{ m}$ والفرق بين مربعي بعديها $16\text{ m}^2$. أجد بعديها. (10 علامات)}
\vspace{7cm}

\newpage
% بداية الصفحة الثالثة
\noindent\textbf{\large السؤال الخامس: إذا كانت النقطة $O$ هي مركز الدائرة في الشكل المجاور. فما قياس الزاويتين المشار إليهما بالرمزين ($x$، $y$)؟ (8 علامات)}
% (ملاحظة: نفس شكل سؤال 1 صفحة 48، الزاوية 30)
\begin{center}
\LTRfig{%
\begin{tikzpicture}[scale=0.9, thick]
  \coordinate (O) at (0,0);
  \coordinate (P) at (180:2.5);   % P على يسار الدائرة
  \coordinate (Q) at (300:2.5);   % Q أسفل اليمين (لأن POQ = 180 - 2*30 = 120)
  \coordinate (R) at (112:2.5);   % رأس الزاوية المحيطية y أعلى اليسار

  \draw[blue!60!black] (O) circle (2.5cm);
  \draw[blue!60!black] (P) -- (O) -- (Q) -- cycle;   % نصفا القطر والوتر PQ
  \draw[blue!60!black] (P) -- (R) -- (Q);            % الزاوية المحيطية

  \fill (O) circle (2pt);
  \node[above right] at (O) {$O$};
  \node[left]        at (P) {$P$};
  \node[below right] at (Q) {$Q$};

  % الزاوية المعطاة عند P
  \pic [draw, red, angle radius=1.1cm, angle eccentricity=1.45, "\small $30^\circ$"] {angle = Q--P--O};
  % الزاوية المركزية x عند O
  \pic [draw, red, angle radius=0.45cm, angle eccentricity=1.7, "$x$"] {angle = P--O--Q};
  % الزاوية المحيطية y عند R
  \pic [draw, red, angle radius=0.6cm, angle eccentricity=1.5, "$y$"] {angle = P--R--Q};
\end{tikzpicture}}
\end{center}
\vspace{2.5cm}

\hrule height 0.4pt
\vspace{0.5cm}

\noindent\textbf{\large السؤال السادس: أسطوانة حجمها $16\pi\text{ cm}^3$، إذا كان طول نصف قطر قاعدة الأسطوانة يقل عن ارتفاعها بـ~$2\text{ cm}$، أجد أبعادها. (4 علامات)}
% (ملاحظة: على نمط سؤال 20 صفحة 15)
\begin{center}
\LTRfig{%
\begin{tikzpicture}[scale=0.9]
    % جسم الأسطوانة (تدرج لوني من الفاتح لليسار إلى الداكن لليمين)
    \shade[left color=cyan!15, right color=teal!45]
        (-2,0) -- (-2,3) arc (180:360:2 and 0.6) -- (2,0) arc (0:-180:2 and 0.6) -- cycle;
    % القاعدة العلوية
    \shade[left color=cyan!25, right color=teal!55] (0,3) ellipse (2 and 0.6);
    \draw[thick, green!50!black] (0,3) ellipse (2 and 0.6);
    % القاعدة السفلية: الجزء الخلفي متقطع والأمامي متصل
    \draw[thick, dashed, green!50!black] (2,0) arc (0:180:2 and 0.6);
    \draw[thick, green!50!black] (-2,0) arc (180:360:2 and 0.6);
    % الجانبان
    \draw[thick, green!50!black] (-2,0) -- (-2,3);
    \draw[thick, green!50!black] (2,0) -- (2,3);
    % نصف القطر r
    \fill (0,0) circle (2pt);
    \draw[thick, red] (0,0) -- (2,0);
    \node[below=3pt] at (1,0) {\large $r$};
    % الارتفاع h
    \draw[thick] (2.5,0) -- (2.5,3) node[midway, right=2pt] {\large $h$};
    \draw[thick] (2.3,0) -- (2.7,0);
    \draw[thick] (2.3,3) -- (2.7,3);
\end{tikzpicture}}
\end{center}
\vspace{3cm}


\vspace{1.5cm}
\begin{center}
\textbf{\large انتهت الأسئلة} \\
\vspace{0.3cm}
\textbf{مع تمنياتي لكنّ بالتوفيق والنجاح} \\
\vspace{0.3cm}
\textbf{إعداد المعلمة: صفا العوضي}
\end{center}

\end{document}
